%==============================================================================
\section{Lựa chọn đặc trưng (Feature Selection)}
\label{sec:feature-selection}
%==============================================================================

\begin{dinhnghia}[title={Lựa chọn đặc trưng}]
Lựa chọn đặc trưng là bước quan trọng trong học máy: chọn một \textbf{tập con} các đặc trưng có sẵn để cải thiện hiệu năng mô hình.
\end{dinhnghia}

\begin{ynghia}[title={Các kỹ thuật phổ biến}]
Dưới đây là một số kỹ thuật lựa chọn đặc trưng thường dùng.
\end{ynghia}

\subsection{Phương pháp lọc (Filter Methods)}

\begin{phuongphap}[title={Ý tưởng Filter}]
Phương pháp này đánh giá mức \textbf{liên quan} của từng đặc trưng bằng cách tính một độ đo thống kê (ví dụ: tương quan, thông tin tương hỗ, kiểm định chi-square, \ldots) rồi \textbf{xếp hạng} đặc trưng theo điểm số. Các đặc trưng có điểm thấp sẽ bị loại khỏi mô hình.
\end{phuongphap}

\begin{vidu}[title={Filter trong Python --- SelectPercentile và chi2}]
Để triển khai phương pháp lọc trong Python, dùng \texttt{SelectKBest} hoặc \texttt{SelectPercentile} từ module \texttt{sklearn.feature\_selection}. Ví dụ chọn 10\% đặc trưng tốt nhất theo chi-square (phù hợp dữ liệu không âm, phân loại):
\begin{lstlisting}
from sklearn.feature_selection import SelectPercentile, chi2

selector = SelectPercentile(chi2, percentile=10)
X_new = selector.fit_transform(X, y)
\end{lstlisting}
\end{vidu}

\begin{luuy}[title={Chi-square và kiểu dữ liệu}]
\texttt{chi2} yêu cầu giá trị đặc trưng \textbf{không âm}. Với đặc trưng liên tục (như bộ Breast Cancer), nên dùng \texttt{f\_classif} hoặc \texttt{mutual\_info\_classif} thay vì \texttt{chi2}.
\end{luuy}

\subsection{Phương pháp bọc (Wrapper Methods)}

\begin{phuongphap}[title={Ý tưởng Wrapper}]
Phương pháp này đánh giá hiệu năng mô hình khi \textbf{thêm hoặc bớt} đặc trưng, rồi chọn tập con đặc trưng cho kết quả tốt nhất. Cách này \textbf{tốn tính toán} hơn phương pháp lọc nhưng thường \textbf{chính xác} hơn.
\end{phuongphap}

\begin{vidu}[title={Wrapper trong Python --- RFE}]
Triển khai wrapper trong Python bằng hàm \texttt{RFE} (Recursive Feature Elimination) trong \texttt{sklearn.feature\_selection}:
\begin{lstlisting}
from sklearn.feature_selection import RFE
from sklearn.linear_model import LogisticRegression

estimator = LogisticRegression()
selector = RFE(estimator, n_features_to_select=5)
selector = selector.fit(X, y)
X_new = selector.transform(X)
\end{lstlisting}
\end{vidu}

\subsection{Phương pháp nhúng (Embedded Methods)}

\begin{phuongphap}[title={Ý tưởng Embedded}]
Phương pháp này \textbf{gắn} lựa chọn đặc trưng vào ngay quá trình xây dựng mô hình, ví dụ hồi quy Lasso, Ridge hoặc cây quyết định. Các phương pháp này gán \textbf{trọng số} cho từng đặc trưng; đặc trưng có trọng số nhỏ sẽ bị loại.
\end{phuongphap}

\begin{vidu}[title={Embedded trong Python --- Lasso}]
Triển khai embedded trong Python bằng \texttt{Lasso} hoặc \texttt{Ridge} từ \texttt{sklearn.linear\_model}:
\begin{lstlisting}
import pandas as pd
from sklearn.linear_model import Lasso

lasso = Lasso(alpha=0.1)
lasso.fit(X, y)
coef = pd.Series(lasso.coef_, index=X.columns)
important_features = coef[coef != 0]
\end{lstlisting}
\end{vidu}

\subsection{Phân tích thành phần chính (PCA)}

\begin{dinhnghia}[title={PCA trong ngữ cảnh lựa chọn đặc trưng}]
PCA là phương pháp học \textbf{không giám sát}: biến đổi các đặc trưng gốc thành tập các thành phần chính \textbf{không tương quan}, giải thích \textbf{phương sai tối đa} trong dữ liệu. Số thành phần chính có thể chọn theo ngưỡng, từ đó \textbf{giảm chiều} tập dữ liệu.
\end{dinhnghia}

\begin{vidu}[title={PCA trong Python}]
Triển khai PCA trong Python bằng \texttt{PCA} từ \texttt{sklearn.decomposition}:
\begin{lstlisting}
from sklearn.decomposition import PCA

pca = PCA(n_components=3)
X_new = pca.fit_transform(X)
\end{lstlisting}
\end{vidu}

\subsection{Loại trừ đặc trưng đệ quy (RFE)}

\begin{phuongphap}[title={RFE --- Recursive Feature Elimination}]
Phương pháp này \textbf{loại dần} các đặc trưng kém quan trọng nhất cho đến khi còn tập đặc trưng quan trọng nhất. RFE dựa trên mô hình; có thể tốn tính toán nhưng cho kết quả tốt trên dữ liệu \textbf{chiều cao}.
\end{phuongphap}

\begin{vidu}[title={RFECV trong Python}]
Triển khai RFE kèm xác thực chéo bằng \texttt{RFECV} (Recursive Feature Elimination with Cross Validation):
\begin{lstlisting}
from sklearn.feature_selection import RFECV
from sklearn.tree import DecisionTreeClassifier

estimator = DecisionTreeClassifier()
selector = RFECV(estimator, step=1, cv=5)
selector = selector.fit(X, y)
X_new = selector.transform(X)
\end{lstlisting}
\end{vidu}

\begin{ynghia}[title={Kết hợp các kỹ thuật}]
Các kỹ thuật lựa chọn đặc trưng trên có thể dùng \textbf{riêng lẻ} hoặc \textbf{kết hợp} để cải thiện hiệu năng mô hình học máy. Cần chọn kỹ thuật phù hợp theo \textbf{kích thước} tập dữ liệu, \textbf{bản chất} đặc trưng và \textbf{loại} mô hình đang dùng.
\end{ynghia}

\subsection{Ví dụ --- Breast Cancer Wisconsin}

\begin{dongco}[title={Bối cảnh ví dụ}]
Ví dụ dưới đây triển khai ba cách: lựa chọn đặc trưng đơn biến (univariate), loại trừ đặc trưng đệ quy có xác thực chéo (RFECV), và phân tích thành phần chính (PCA).
\end{dongco}

\begin{ynghia}[title={Bộ dữ liệu Breast Cancer}]
Dùng bộ \textbf{Breast Cancer Wisconsin (Diagnostic)} có sẵn trong scikit-learn: 569 mẫu, 30 đặc trưng; nhiệm vụ phân loại khối u \textbf{ác tính} (malignant) hay \textbf{lành tính} (benign).
\end{ynghia}

\begin{vidu}[title={Mã Python đầy đủ}]
\begin{lstlisting}
# Import thu vien va bo du lieu
import numpy as np
from sklearn.datasets import load_breast_cancer
from sklearn.feature_selection import SelectKBest, f_classif, RFECV
from sklearn.model_selection import train_test_split, cross_val_score
from sklearn.linear_model import LogisticRegression
from sklearn.decomposition import PCA

# Tai bo du lieu
data = load_breast_cancer()
X = data.data
y = data.target

# Loc don bien: f_classif (phu hop dac trung lien tuc)
selector = SelectKBest(f_classif, k=10)
X_new = selector.fit_transform(X, y)

# Chia tap huan luyen / kiem tra
X_train, X_test, y_train, y_test = train_test_split(
    X_new, y, test_size=0.3, random_state=42)

# Hoi quy logistic tren dac trung da chon
clf = LogisticRegression(max_iter=10000)
clf.fit(X_train, y_train)

# Danh gia tren tap kiem tra
accuracy = clf.score(X_test, y_test)
print("Accuracy using univariate feature selection: {:.2f}".format(accuracy))

# RFECV
estimator = LogisticRegression(max_iter=10000)
selector = RFECV(estimator, step=1, cv=5)
selector.fit(X, y)
X_new = selector.transform(X)
scores = cross_val_score(LogisticRegression(max_iter=10000), X_new, y, cv=5)
print("Accuracy using RFECV feature selection: %0.2f (+/- %0.2f)" %
      (scores.mean(), scores.std() * 2))

# PCA
pca = PCA(n_components=5)
X_new = pca.fit_transform(X)
scores = cross_val_score(LogisticRegression(max_iter=10000), X_new, y, cv=5)
print("Accuracy using PCA feature selection: %0.2f (+/- %0.2f)" %
      (scores.mean(), scores.std() * 2))
\end{lstlisting}
\end{vidu}

\begin{luuy}[title={Sua loi so voi ban goc Tutorialspoint}]
Ban goc dung \texttt{chi2} va duong dan \texttt{C:\textbackslash...} tren file diabetes.csv; day la \textbf{Breast Cancer} voi dac trung \textbf{lien tuc}, nen dung \texttt{load\_breast\_cancer} va \texttt{f\_classif}. Them import \texttt{cross\_val\_score} va \texttt{PCA} con thieu trong ban goc.
\end{luuy}

\begin{vidu}[title={Ket qua chay (tham khao)}]
Khi chay ma tren, terminal co the in tuong tu:
\begin{lstlisting}
Accuracy using univariate feature selection: 0.74
Accuracy using RFECV feature selection: 0.77 (+/- 0.03)
Accuracy using PCA feature selection: 0.75 (+/- 0.07)
\end{lstlisting}
\end{vidu}

\begin{tinhchat}[title={So sanh nhanh ba cach trong vi du}]
\begin{itemize}
  \item \textbf{Univariate + f\_classif}: nhanh, doc lap mô hình; chon $k=10$ trong 30 dac trung.
  \item \textbf{RFECV}: tim tap dac trung tot theo hieu nang mo hinh va CV; ton tinh toan hon.
  \item \textbf{PCA}: giam chieu bang bien doi tuyen tinh; khong giu ten dac trung goc.
\end{itemize}
\end{tinhchat}
